% Seminar handout template (1–2 pages) — copy to NN-topic/handout.tex.
\documentclass{ca-handout}

\seminar{0}{Быстрое возведение в степень}

\begin{document}
\makehandouttitle

\section{Бинарный алгоритм}
Пусть $n = \sum_{i=0}^{k} n_i 2^i$, $n_i \in \{0, 1\}$. Тогда
\[
  a^n = \prod_{i\,:\,n_i = 1} a^{2^i}.
\]

\begin{reminder}[title={Алгоритм «справа налево»}]
  $r \gets 1$; для $i = 0, \dots, k$: если $n_i = 1$, то $r \gets r \cdot a$; $a \gets a^2$.
  Сложность: $\bigO(\log n)$ умножений.
\end{reminder}

\section{Метод Монтгомери}
Для нечётного $m$ и $R = 2^k > m$, $\gcd(R, m) = 1$:
\[
  \operatorname{REDC}(T) = T R^{-1} \bmod m, \qquad 0 \le T < mR.
\]

\checkyourself{Почему для метода Монтгомери модуль $m$ должен быть нечётным?}
\end{document}
